\documentclass[11pt]{article}

% ============================================================
% COMPILATION
% ============================================================
%
% Required compiler: XeLaTeX
%
% Recommended build sequence:
% XeLaTeX -> Biber -> XeLaTeX -> XeLaTeX
%

% ============================================================
% DOCUMENT EDITION
% ============================================================

\newif\iftutornotes

\tutornotestrue       % TA / Tutor edition
% \tutornotesfalse    % Student edition

% ============================================================
% PACKAGES
% ============================================================

\usepackage[a4paper,margin=1in]{geometry}

% XeLaTeX font system
\usepackage{fontspec}
\usepackage{unicode-math}

\setmainfont{TeX Gyre Termes}
\setsansfont{TeX Gyre Heros}
\setmonofont{Latin Modern Mono}
\setmathfont{TeX Gyre Termes Math}

\usepackage{microtype}
\usepackage{enumitem}
\usepackage{booktabs}
\usepackage{array}
\usepackage{longtable}
\usepackage{xcolor}
\usepackage{hyperref}
\usepackage{listings}
\usepackage{tcolorbox}
\usepackage{titlesec}
\usepackage{fancyhdr}
\usepackage[backend=biber,style=ieee]{biblatex}

% ============================================================
% DOCUMENT METADATA
% ============================================================

\newcommand{\CourseTitle}{Physical AI, TAICA}
\newcommand{\InstructorName}{Human-centered Intelligent System Lab}
\newcommand{\InstructorEmail}{National Yang Ming Chiao Tung University}

\newcommand{\AssignmentTitle}{Homework N: Assignment Title}
\newcommand{\Semester}{Fall 2026}

% ============================================================
% BIBLIOGRAPHY
% ============================================================

\addbibresource{refs.bib}

% ============================================================
% TABLE SETTINGS
% ============================================================

\newcolumntype{L}[1]{>{\raggedright\arraybackslash}p{#1}}

% ============================================================
% HYPERLINK SETTINGS
% ============================================================

\hypersetup{
    colorlinks=true,
    linkcolor=blue!50!black,
    citecolor=blue!50!black,
    urlcolor=blue!50!black
}

% ============================================================
% PAGE STYLE
% ============================================================

\setlength{\headheight}{14pt}

\pagestyle{fancy}
\fancyhf{}

\lhead{\CourseTitle}
\rhead{\AssignmentTitle}
\cfoot{\thepage}

% ============================================================
% LIST SETTINGS
% ============================================================

\setlist[itemize]{
    topsep=3pt,
    itemsep=2pt,
    parsep=0pt,
    leftmargin=1.5em
}

\setlist[enumerate]{
    topsep=3pt,
    itemsep=2pt,
    parsep=0pt,
    leftmargin=1.6em
}

% ============================================================
% COLORS
% ============================================================

\definecolor{codebg}{RGB}{248,248,248}
\definecolor{commentgreen}{RGB}{0,110,60}
\definecolor{keywordblue}{RGB}{20,70,150}

% A single restrained color for all internal teaching notes.
\definecolor{tutorblue}{RGB}{65,88,110}
\definecolor{tutorbg}{RGB}{247,249,251}

% ============================================================
% TA / TUTOR NOTE SYSTEM
% ============================================================

% All material intended only for teaching staff must use this
% environment. Do not use separate teacher/grading/TA colors.
%
% In student mode, the complete environment body disappears.

\NewDocumentEnvironment{tutornote}{+b}
{%
    \iftutornotes
        \begin{tcolorbox}[
            colback=tutorbg,
            colframe=tutorblue,
            boxrule=0.4pt,
            arc=1mm,
            left=7pt,
            right=7pt,
            top=6pt,
            bottom=6pt,
            before skip=8pt,
            after skip=8pt
        ]
        {\small
        \color{tutorblue}
        \textbf{TA/Tutor Note.}\par
        \medskip
        #1}
        \end{tcolorbox}
    \fi
}
{}

% Optional edition indicator.
% It appears only in the teaching edition.

\newcommand{\EditionLabel}{%
    \iftutornotes
        \begin{center}
            {\small\color{tutorblue}
            TA / Tutor Teaching Edition}
        \end{center}
        \vspace{0.5em}
    \fi
}

% ============================================================
% LISTINGS
% ============================================================

\lstdefinelanguage{turtle}{
    alsoletter={:@},
    morekeywords={
        @prefix,
        a,
        owl:,
        rdf:,
        rdfs:,
        xsd:
    },
    sensitive=true,
    morecomment=[l]{\#},
    morestring=[b]"
}

\lstset{
    basicstyle=\ttfamily\small,
    backgroundcolor=\color{codebg},
    breaklines=true,
    frame=single,
    framerule=0.2pt,
    rulecolor=\color{black!20},
    keywordstyle=\color{keywordblue}\bfseries,
    commentstyle=\color{commentgreen},
    stringstyle=\color{black!70},
    columns=fullflexible,
    keepspaces=true,
    showstringspaces=false
}

\lstdefinestyle{codesnippet}{
    numbers=none,
    aboveskip=0.75\baselineskip,
    belowskip=0.85\baselineskip,
    frame=single,
    framerule=0.2pt,
    rulecolor=\color{black!20},
    backgroundcolor=\color{codebg},
    basicstyle=\ttfamily\small,
    breaklines=true,
    columns=fullflexible,
    keepspaces=true,
    showstringspaces=false
}

\newcommand{\code}[1]{\texttt{\detokenize{#1}}}

% ============================================================
% TITLE
% ============================================================

\title{\textbf{\AssignmentTitle}}
\author{
    \InstructorName\\
    \InstructorEmail
}
\date{\Semester}

% ============================================================
% DOCUMENT
% ============================================================

\begin{document}

\maketitle
\EditionLabel

% ============================================================
\section*{Introduction}
% ============================================================

This assignment introduces a selected topic through a combination of
conceptual analysis, technical implementation, and documented results.
Students are expected to produce a small but complete artifact that
demonstrates the concepts introduced in class.

The exact subject matter, required resources, and submission format should
be adapted to the objectives of the course.

\begin{tutornote}
Use this space for teaching information that should not appear in the
student handout.

Typical content may include the pedagogical purpose of the assignment,
connections to previous lectures, concepts that require additional
explanation, or suggestions for conducting the tutorial session.

The student-facing text should remain self-contained even when all such
notes are removed.
$$
x+y=z
$$
\end{tutornote}

% ============================================================
\section{Context and Motivation}
% ============================================================

Explain the disciplinary or technical context of the assignment and why the
problem matters.

A useful conceptual structure is:

\[
\text{course concept}
\longrightarrow
\text{practical problem}
\longrightarrow
\text{student exercise}.
\]

Students should be able to understand how the assignment fits into the wider
course rather than treating it as an isolated technical exercise.

\begin{tutornote}
During the tutorial, connect this section explicitly to the relevant lecture
material. If students have difficulty identifying the purpose of the task,
ask them first to distinguish the general course problem from the particular
artifact they are being asked to construct.
\end{tutornote}

% ============================================================
\section{Learning Objectives}
% ============================================================

After completing this assignment, students should be able to:

\begin{enumerate}
    \item explain the principal concepts introduced in the assignment;
    \item distinguish several related conceptual categories;
    \item construct a small technical or analytical artifact;
    \item evaluate the resulting artifact using an explicit procedure;
    \item document assumptions, methods, and results reproducibly;
    \item discuss limitations and possible extensions.
\end{enumerate}

\begin{tutornote}
When adapting this template, formulate learning objectives using observable
actions such as \emph{identify}, \emph{construct}, \emph{compare},
\emph{evaluate}, or \emph{justify}. These objectives should correspond to
what is actually assessed later in the rubric.
\end{tutornote}

% ============================================================
\section{Conceptual Foundation}
% ============================================================

\subsection*{Concept A}

Introduce the first concept required for the assignment. Provide definitions,
notation, examples, and references where appropriate
\cite{example_reference}.

\subsection*{Concept B}

Introduce a second concept and explain how it differs from Concept A.
The distinction should be relevant to the work students will perform.

For example:

\[
A \neq B,
\]

even if the two concepts frequently occur together in practical systems.

\begin{tutornote}
This is an appropriate location for a fuller conceptual explanation that a
TA or tutor may need when answering questions.

It may contain several paragraphs, equations, or lists:

\begin{itemize}
    \item explain the intended distinction;
    \item identify common misconceptions;
    \item indicate what level of formal detail is expected;
    \item record examples that may be used during discussion.
\end{itemize}
\end{tutornote}

% ============================================================
\section{Problem Statement}
% ============================================================

Assume a system containing several entities with observable attributes and
relationships.

Your task is to construct a model that represents the relevant entities,
processes them according to the methods introduced in class, and produces
results that can be independently inspected.

Your solution must:

\begin{enumerate}
    \item represent the required entities;
    \item implement at least one non-trivial relation or operation;
    \item produce a reproducible result;
    \item demonstrate the result using an appropriate query, test, or analysis;
    \item explain the principal design decisions.
\end{enumerate}

\begin{tutornote}
The distinction between ``non-trivial result'' and manually supplied output
should be made explicit during the tutorial.

A student should be able to explain which input facts, rules, transformations,
or computations are responsible for the reported result.
\end{tutornote}

% ============================================================
\section{Required Resources}
% ============================================================

The submitted work must contain, at minimum:

\begin{itemize}
    \item one primary model or data file;
    \item human-readable names or annotations for major entities;
    \item at least one explicit structural relation;
    \item at least one result that is derived rather than manually copied;
    \item at least one query, test, analysis, or validation procedure;
    \item sufficient documentation for reproduction.
\end{itemize}

% ============================================================
\section{Recommended Modeling Distinctions}
% ============================================================

A useful design principle is to distinguish the different kinds of entities
represented in the system.

{\renewcommand{\arraystretch}{1.25}
\begin{longtable}{L{0.20\linewidth}L{0.73\linewidth}}
\toprule
\textbf{Layer} & \textbf{Meaning} \\
\midrule

Type &
A general category shared by multiple entities. \\

Role &
A context-dependent function that an entity performs in a task. \\

Property &
An attribute or relation describing an entity. \\

Instance &
A particular entity occurring in the dataset or environment. \\

Derived result &
A value or classification obtained through inference, computation, or
analysis rather than directly supplied as input. \\

\bottomrule
\end{longtable}
}

\begin{tutornote}
Replace these categories with distinctions appropriate to the assignment.

The purpose of the table is pedagogical: it should expose distinctions that
students are otherwise likely to collapse into a single informal category.
\end{tutornote}

% ============================================================
\section{Base Example}
% ============================================================

Students may use the following generic structure as a starting point:

\begin{lstlisting}[
    caption={Generic structured-data example.},
    label={lst:generic-example}
]
# Example only

example01:
    type: ExampleType
    label: "Example entity"
    value: 42
    related_to: example02
\end{lstlisting}

The example is deliberately incomplete. Students should extend or replace it
according to the requirements of the assignment.

\begin{tutornote}
Clarify how much modification is necessary for the example to constitute a
student-authored solution. If merely copying the example would satisfy the
formal requirements, the example is probably too complete.
\end{tutornote}

% ============================================================
\section{Identifier or Namespace Policy}
% ============================================================

If the assignment uses identifiers, namespaces, package names, database
schemas, URI patterns, or other globally scoped identifiers, define the
convention here.

For example:

\begin{lstlisting}[style=codesnippet]
shared:   https://example.edu/course/shared/
groupNN:  https://example.edu/course/groupNN/
\end{lstlisting}

Shared resources should use the shared namespace. Student-created entities
should use a group-specific namespace.

% ============================================================
\section{Input Data or Environment}
% ============================================================

Students must provide sufficient input data to demonstrate the required
functionality.

Input data should normally be stored separately from generated results.

A minimal example might be:

\begin{lstlisting}[style=codesnippet]
item01:
    type: Example
    label: "Example item"
    value: 42
\end{lstlisting}

% ============================================================
\section{Required Processing or Reasoning Pattern}
% ============================================================

The assignment should contain at least one relationship whose result is not
simply asserted manually. Students should be able to explain the same modeling
commitment at the conceptual, formal, and implementation levels.

Suppose the assignment introduces the following general pattern:

\begin{quote}
An entity is classified as a \emph{DerivedEntity} if it is a
\emph{BaseEntity} and participates in at least one
\emph{relevantRelation} with an entity of type \emph{RelatedEntity}.
\end{quote}

At the predicate-logic level, this may be written as:

\[
\mathrm{DerivedEntity}(x)
\iff
\mathrm{BaseEntity}(x)
\land
\exists y\,
\bigl(
    \mathrm{relevantRelation}(x,y)
    \land
    \mathrm{RelatedEntity}(y)
\bigr).
\]

The same condition can be represented more compactly in Description Logic as:

\[
\boxed{
\mathsf{DerivedEntity}
\equiv
\mathsf{BaseEntity}
\sqcap
\exists\,\mathsf{relevantRelation}.\mathsf{RelatedEntity}
}
\]

Here, $\sqcap$ denotes class intersection, while $\exists\,R.C$ denotes an existential restriction: the set of individuals that participate
in at least one relation \(R\) to an individual belonging to class \(C\).

Thus,
\[
\exists\,\mathsf{relevantRelation}.\mathsf{RelatedEntity}
\]

denotes the class of all individuals that have at least one
\code{ex:relevantRelation} relation to an instance of
\code{ex:RelatedEntity}.

The corresponding OWL 2 class definition can be serialized in Turtle as
follows:

\begin{lstlisting}[
    language=turtle,
    caption={OWL/Turtle representation of a Description Logic class definition.},
    label={lst:dl-owl-example}
]
@prefix ex:   <https://example.org/ontology/> .
@prefix owl:  <http://www.w3.org/2002/07/owl#> .
@prefix rdfs: <http://www.w3.org/2000/01/rdf-schema#> .

ex:DerivedEntity
    a owl:Class ;
    rdfs:label "derived entity"@en ;
    owl:equivalentClass [
        a owl:Class ;
        owl:intersectionOf (
            ex:BaseEntity
            [
                a owl:Restriction ;
                owl:onProperty ex:relevantRelation ;
                owl:someValuesFrom ex:RelatedEntity
            ]
        )
    ] .
\end{lstlisting}

The three representations therefore correspond as follows:

{\renewcommand{\arraystretch}{1.3}
\begin{longtable}{L{0.23\linewidth}L{0.70\linewidth}}
\toprule
\textbf{Representation} & \textbf{Expression} \\
\midrule

Informal &
An entity is a DerivedEntity if it is a BaseEntity and has at least one
relevantRelation to a RelatedEntity. \\

Predicate logic &
\(
\mathrm{DerivedEntity}(x)
\iff
\mathrm{BaseEntity}(x)
\land
\exists y\,
(
\mathrm{relevantRelation}(x,y)
\land
\mathrm{RelatedEntity}(y)
)
\) \\

Description Logic &
\(
\mathsf{DerivedEntity}
\equiv
\mathsf{BaseEntity}
\sqcap
\exists\,\mathsf{relevantRelation}.\mathsf{RelatedEntity}
\) \\

OWL 2 &
\code{owl:equivalentClass},
\code{owl:intersectionOf},
\code{owl:Restriction},
\code{owl:onProperty}, and
\code{owl:someValuesFrom}. \\

\bottomrule
\end{longtable}
}

The intended inference is therefore:

\[
\mathsf{BaseEntity}(a)
\quad+\quad
\mathsf{relevantRelation}(a,b)
\quad+\quad
\mathsf{RelatedEntity}(b)
\]
\[
\Downarrow
\]
\[
\boxed{
\mathsf{DerivedEntity}(a)
}
\]

For example, the asserted instance data may be:

\begin{lstlisting}[
    language=turtle,
    caption={Example asserted individuals.}
]
ex:item01
    a ex:BaseEntity ;
    ex:relevantRelation ex:relatedItem01 .

ex:relatedItem01
    a ex:RelatedEntity .
\end{lstlisting}

If a reasoner supports the OWL constructs used above, it may therefore infer:

\begin{lstlisting}[
    language=turtle,
    caption={Example inferred class membership.}
]
ex:item01
    a ex:DerivedEntity .
\end{lstlisting}

Students should distinguish the final inferred triple from the asserted input.
The point of the exercise is not merely to write the expected classification
into the RDF graph, but to define formal axioms from which the classification
can be derived.

\begin{tutornote}
This example is intended to demonstrate the correspondence between different
levels of formalization.

The important conceptual mapping is:

\[
\begin{aligned}
\text{class conjunction}
&\longleftrightarrow
\sqcap
\longleftrightarrow
\code{owl:intersectionOf},
\\[4pt]
\text{existential restriction}
&\longleftrightarrow
\exists R.C
\longleftrightarrow
\code{owl:someValuesFrom},
\\[4pt]
\text{class equivalence}
&\longleftrightarrow
\equiv
\longleftrightarrow
\code{owl:equivalentClass}.
\end{aligned}
\]

During teaching, emphasize that Turtle is only a serialization syntax. The
formal meaning is supplied by the OWL semantics, while Description Logic gives
a compact notation for the same class expression.

Students should therefore not infer that the Turtle syntax itself constitutes
the reasoning formalism.
\end{tutornote}

% ============================================================
\section{Query or Evaluation Requirement}
% ============================================================

Each submission must include at least one query, test, or evaluation procedure
that demonstrates the required result.

For example:

\begin{lstlisting}[
    language=SQL,
    caption={Generic query example.}
]
SELECT item, result
FROM derived_results
WHERE result IS NOT NULL;
\end{lstlisting}

Expected output should contain at least one result that can be traced back to
the submitted model and processing procedure.

% ============================================================
\section{Technology Options}
% ============================================================

Students may use one of the supported workflows.

\subsection{Option A: Graphical Tool}

\begin{enumerate}
    \item Open or create the project.
    \item Load the required input.
    \item Inspect the model or data structure.
    \item Execute the required processing step.
    \item Verify and export the result.
\end{enumerate}

\subsection{Option B: Command-Line Workflow}

A command-line workflow may resemble:

\begin{lstlisting}[language=bash]
tool --validate input/example.dat

tool \
    --input input/example.dat \
    --output results/output.dat
\end{lstlisting}

\subsection{Option C: Programming-Language Workflow}

Students may implement the workflow in a supported programming language,
provided that dependencies and execution commands are documented clearly.

\begin{tutornote}
Before the assignment session, verify that the listed tools remain available
and that the example commands still work with the versions used in the
teaching environment.

Keep implementation-specific troubleshooting information here rather than
placing it in the student instructions unless students genuinely need it.
\end{tutornote}

% ============================================================
\section{Minimal Programming Example}
% ============================================================

The following skeleton illustrates a generic processing workflow:

\begin{lstlisting}[
    language=Python,
    caption={Minimal processing skeleton.}
]
def load_input(path):
    # Load student data.
    pass

def process(data):
    # Apply the required transformation.
    pass

def evaluate(result):
    # Produce the required output.
    pass

def main():
    data = load_input("input/example.dat")
    result = process(data)
    evaluate(result)

if __name__ == "__main__":
    main()
\end{lstlisting}

% ============================================================
\section{Optional Validation}
% ============================================================

Additional validation may be included where appropriate.

Possible checks include:

\begin{itemize}
    \item required fields are present;
    \item identifiers are unique;
    \item values satisfy their intended type constraints;
    \item relationships point to valid entities;
    \item generated results can be reproduced from the submitted input.
\end{itemize}

% ============================================================
\section{Deliverables}
% ============================================================

Each group must submit a repository containing the artifacts required to
reproduce the work.

\subsection{Repository Structure}

\begin{lstlisting}[language=bash]
assignment-project/
|-- README.md
|-- report.pdf
|-- input/
|   `-- example.dat
|-- model/
|   `-- model.ext
|-- queries/
|   `-- example-query.ext
|-- results/
|   `-- output.txt
|-- src/
|   `-- ...
`-- LICENSE
\end{lstlisting}

\subsection{README Requirements}

The \code{README.md} should include:

\begin{enumerate}
    \item project title and group members;
    \item selected problem or task;
    \item design overview;
    \item description of major files;
    \item execution instructions;
    \item expected output;
    \item explanation of derived results;
    \item known limitations.
\end{enumerate}

\subsection{Required Submission Components}

{\renewcommand{\arraystretch}{1.25}
\begin{longtable}{L{0.25\linewidth}L{0.68\linewidth}}
\toprule
\textbf{Repository item} & \textbf{Requirement} \\
\midrule

\code{README.md} &
Explains the project, repository structure, execution procedure, and expected
results. \\

\code{report.pdf} &
Explains the design rationale, implementation, results, limitations, and
relevant discussion. \\

\code{model/} &
Contains the principal student-authored representation. \\

\code{queries/} &
Contains the required queries, tests, or analysis procedures. \\

\code{results/} &
Contains saved output sufficient to demonstrate the reported results. \\

\code{src/} &
Contains source code or configuration files when applicable. \\

\bottomrule
\end{longtable}
}

% ============================================================
\section{Assessment Rubric}
% ============================================================

\begin{longtable}{
    p{0.25\linewidth}
    p{0.10\linewidth}
    p{0.55\linewidth}
}
\toprule

\textbf{Criterion}
&
\textbf{Weight}
&
\textbf{Indicators}
\\

\midrule

Completeness
&
25\%
&
Required resources are present and the submitted artifact covers the
specified problem. \\

Conceptual quality
&
25\%
&
Important categories and relationships are distinguished correctly and
meaningfully. \\

Processing and evaluation
&
30\%
&
At least one non-trivial result is derived and demonstrated reproducibly. \\

Documentation
&
20\%
&
Repository structure, README, report, and execution instructions are complete
and understandable. \\

\bottomrule
\end{longtable}

\begin{tutornote}
The rubric above is visible to students and should therefore contain the
official assessment criteria.

Use this note only for internal interpretation of those criteria: borderline
cases, expected evidence, recurrent grading issues, or examples of what should
and should not receive full credit.

Do not introduce a hidden criterion here that is absent from the published
rubric.
\end{tutornote}

% ============================================================
\section{Common Pitfalls}
% ============================================================

\begin{enumerate}

    \item
    \textbf{Manually inserting expected results.}
    Results that are supposed to be derived should not simply be copied into
    the output.

    \item
    \textbf{Confusing categories and individual examples.}
    A general concept and a particular instance should be represented
    separately.

    \item
    \textbf{Using labels as a substitute for structure.}
    Human-readable text does not replace explicit relationships in the model.

    \item
    \textbf{Reporting results without a reproducible procedure.}
    The submitted files should make it possible to regenerate the result.

\end{enumerate}

\begin{tutornote}
This section can accumulate observations from successive offerings of the
course.

For example, record which mistakes are genuinely conceptual and which usually
result from software configuration. This information can help tutors diagnose
student problems without adding unnecessary detail to the student handout.
\end{tutornote}

% ============================================================
\section{Suggested Extensions}
% ============================================================

Students seeking an advanced implementation may:

\begin{itemize}
    \item introduce additional entity or relation types;
    \item implement stronger validation;
    \item compare alternative modeling approaches;
    \item automate the processing pipeline;
    \item evaluate the model against additional data;
    \item discuss scalability or generalization.
\end{itemize}

% ============================================================
\section{Conclusion}
% ============================================================

This assignment connects conceptual understanding with a reproducible technical
artifact.

A successful submission should make the relationship between

\[
\text{input}
\rightarrow
\text{representation}
\rightarrow
\text{processing}
\rightarrow
\text{result}
\rightarrow
\text{evaluation}
\]

explicit and independently inspectable.

% ============================================================
% APPENDIX
% ============================================================

\appendix

\section{Additional Reference Material}

Optional material that is useful for completing the assignment but would
interrupt the main exposition may be placed in the appendix.

\begin{tutornote}
Teaching-only appendix material can also be kept here, including:

\begin{itemize}
    \item expected solution structure;
    \item worked examples;
    \item tutorial prompts;
    \item likely student questions;
    \item notes inherited from previous course offerings.
\end{itemize}

Because the complete environment disappears when
\code{\string\tutornotesfalse} is selected, the same source file can remain the
master copy for both editions.
\end{tutornote}

% ============================================================
% BIBLIOGRAPHY
% ============================================================

\printbibliography

\end{document}